\documentclass[runningheads]{llncs}

% ECCVW AI City Challenge submissions are named rather than double-blind. The title, abstract,
% keywords, author names, and affiliations are immutable after initial submission. Keep the main
% manuscript within the 14-page limit; references and supplementary material follow workshop rules.
\usepackage{eccv}
\usepackage{eccvabbrv}
\usepackage{graphicx}
\usepackage{booktabs}
\usepackage{amsmath}
\usepackage{xcolor}
\usepackage{hyperref}

\newcommand{\todo}[1]{\textcolor{red}{[TBD: #1]}}

\begin{document}

\title{STAR: Action-Aware Two-Stage Retrieval for Fine-Grained Person Anomaly Retrieval}
\titlerunning{STAR for Fine-Grained Person Anomaly Retrieval}

\author{Tuan Vo-Lan \and Khanh Tran-Quoc \and Phuc Nguyen-Ngoc \and Thanh Vo-Thai}
\authorrunning{Tuan Vo-Lan et al.}
\institute{University of Science, Vietnam National University Ho Chi Minh City}

\maketitle

\begin{abstract}
Fine-grained person anomaly retrieval requires matching a textual description to an image whose
decisive evidence may be an action, interaction, or event state rather than its global scene.
We present STAR, an action-aware two-stage retrieval framework. STAR enriches anomaly-side
training captions with source event information, adapts a PE-Core visual retriever with
Cross-Batch Memory for high-recall candidate generation, and reranks candidates with an X-VLM
cross-encoder optimized by contrastive, image-text matching, and Smooth-AP objectives.
A Gale--Shapley assignment stage resolves conflicting Top-1 image choices under the one-to-one
retrieval structure. On a clean video-disjoint held-out split, STAR obtains \todo{metrics}.
\keywords{Text-image retrieval \and person anomaly retrieval \and hard negatives \and global assignment}
\end{abstract}

\section{Introduction}
Text-based person anomaly retrieval is difficult when candidates share a scene, subjects,
objects, and background while differing only in action phase, body state, object contact, or
event outcome. Global contrastive retrieval can locate the correct neighborhood but frequently
confuses near-identical candidates. STAR addresses this issue with action-aware captions,
high-recall retrieval, cross-encoder verification, and a final one-to-one assignment.

Our contributions are: (i) an action-aware PE-to-X-VLM retrieval architecture optimized with
contrastive, image-text matching, and rank-aware objectives; (ii) anomaly-focused caption
enhancement using source event descriptions; (iii) PE adaptation with Cross-Batch Memory; and
(iv) a Gale--Shapley postprocessor for one-to-one Top-1 resolution.

\section{Related Work}
\paragraph{Text-image retrieval.} Discuss dual encoders, cross-encoders, and the efficiency--
accuracy trade-off motivating two-stage retrieval. \todo{Add citations.}
\paragraph{Fine-grained retrieval and ranking.} Discuss hard negatives and differentiable ranking,
including Smooth-AP. \todo{Add citations.}
\paragraph{Caption augmentation and global assignment.} Discuss event-aware language supervision,
similarity-coverage analysis, and stable matching. \todo{Add citations.}

\section{Method}
\subsection{Data Preparation and Caption Enhancement}
We construct deterministic video-disjoint train/dev/report partitions. All hard-pair endpoints
remain within one partition. Nested 10K, 30K, and 50K hard subsets are sampled exclusively from
training videos. Anomaly-side captions are rewritten from the original caption and source event
description; normal captions are preserved unchanged.

\subsection{Overall Architecture}
STAR first encodes images and captions with PE-Core to obtain a high-recall candidate set. The
X-VLM cross-encoder reranks the Top-$K$ candidates using fine-grained image-text compatibility.
Gale--Shapley then resolves duplicate Top-1 assignments while retaining the score ordering of
the remaining candidates. The first stage is recall-oriented, whereas the second stage is trained
to resolve the action and state distinctions within a small candidate set. Figure~\ref{fig:pipeline}
will visualise the pipeline.

\begin{figure}[t]
  \centering
  \fbox{\rule{0pt}{1.1in}\rule{0.95\linewidth}{0pt}}
  \caption{\todo{Pipeline figure: caption enhancement $\rightarrow$ PE + XBM $\rightarrow$
  Top-$K$ $\rightarrow$ X-VLM ITM $\rightarrow$ Gale--Shapley.}}
  \label{fig:pipeline}
\end{figure}

\subsubsection{PE First-Stage Retrieval}
PE-Core-bigG-14-448 produces normalized image and caption embeddings. We freeze and precompute
the text representation, and adapt the visual tower, projection layers, and visual LoRA modules.
With Cross-Batch Memory, its symmetric contrastive objective is
\begin{equation}
\mathcal{L}_{\mathrm{PE}}=\tfrac{1}{2}\left[
\mathrm{CE}(S_{I\rightarrow T}, y)+\mathrm{CE}(S_{T\rightarrow I}, y)\right].
\end{equation}
This stage returns Top-$K$ candidates rather than a final one-to-one prediction.

\subsubsection{X-VLM Cross-Encoder and Explicit Hard Pairs}
The X-VLM cross-encoder scores candidate image-text pairs with an ITM head. For every matched
pair $(I_i,T_i)$ and its linked hard partner $j$, we train on the positive $(I_i,T_i)$ and the two
directed negatives $(I_j,T_i)$ and $(I_i,T_j)$. This targets action and state ambiguity within
otherwise similar scenes.

\subsubsection{Optimization Objectives}
The X-VLM stage minimizes
\begin{equation}
\mathcal{L}=\mathcal{L}_{\mathrm{ITC}}+
\lambda_{\mathrm{ITM}}\mathcal{L}_{\mathrm{ITM}}+
\lambda_{\mathrm{SAP}}\mathcal{L}_{\mathrm{SmoothAP}}.
\end{equation}
ITC aligns global image and text embeddings, ITM compares positives with explicit hard image and
hard text negatives, and Smooth-AP supplies a differentiable early-rank signal. All loss weights
are selected on dev only. We do not use masked language modeling during PE or X-VLM fine-tuning:
the text encoder is frozen and the X-VLM objective contains exactly the three terms above.

\subsection{Cross-Batch Memory}
PE adaptation maintains FIFO queues of detached normalized image and text features. Instance and
caption identifiers mask invalid same-instance and duplicate-caption negatives, enlarging the
contrastive negative set beyond the physical batch.

\subsection{Gale--Shapley Assignment}
Each query proposes to candidates in descending reranker-score order. An image holds its strongest
proposal and rejects weaker ones until convergence. The held candidate becomes rank one; the
remaining list preserves score order after that candidate is removed.

\section{Experiments}
\subsection{Protocol and Metrics}
All model-selection decisions use dev. The report split is evaluated only after a recipe is frozen.
We report mAP, R@1, R@5, and R@10.

\subsection{Experimental Settings}
\todo{Fill from frozen configs: PE 448px, X-VLM 384px, LoRA, augmentation, XBM, batch,
optimizer, schedule, seed, and compute.}

\subsection{Results and Ablations}
\begin{table}[t]
  \centering
  \caption{Main comparison on the frozen report split.}
  \begin{tabular}{lrrrr}
    \toprule
    Method & mAP & R@1 & R@5 & R@10 \\
    \midrule
    PE retrieval & \todo{} & \todo{} & \todo{} & \todo{} \\
    + X-VLM reranking & \todo{} & \todo{} & \todo{} & \todo{} \\
    + Gale--Shapley (STAR) & \todo{} & \todo{} & \todo{} & \todo{} \\
    \bottomrule
  \end{tabular}
  \label{tab:main}
\end{table}

\todo{Add scale, component-ablation, and postprocessing tables from the experiment ledger.}

\section{Qualitative Analysis and Limitations}
\todo{Add report-split retrieval wins and two failure cases. Discuss ambiguous event phase,
caption ambiguity, and domain shift.}

\bibliographystyle{splncs04}
\bibliography{main}
\end{document}
